\documentclass[../main]{subfiles}
\begin{document}

\chapter{Teorema de Bernoulli}
\section{Teoría}

Teorema de Bernoulli
El Teorema de Bernoulli es uno de los principios fundamentales de la mecánica de fluidos. Establece una relación entre la velocidad, la presión y la altura de un fluido en movimiento. Este teorema es una manifestación del principio de conservación de la energía y es especialmente útil en el análisis de flujos incomprensibles y no viscosos.

El Teorema de Bernoulli se expresa matemáticamente como:
\begin{equation}
P + \frac{1}{2} \rho v^2 + \rho g h = \text{constante}
\end{equation}

Donde:
\begin{itemize}
    \item $P$: Presión del fluido [Pa]
    \item $\rho$: Densidad del fluido [kg/m^3]
    \item $v$: Velocidad del fluido [m/s]
    \item $g$: Aceleración debida a la gravedad [9.81 m/s^2]
    \item $h$: Altura sobre un punto de referencia [m]
\end{itemize}

Para entender mejor cómo se aplica el Teorema de Bernoulli, consideremos un fluido que fluye a través de una tubería de diferentes diámetros. En la sección de la tubería donde el diámetro es menor, la velocidad del fluido aumenta y la presión disminuye, y viceversa. Este comportamiento se puede analizar con la ecuación de continuidad y la ecuación de Bernoulli.

La ecuación de continuidad para un fluido incompresible es:
\begin{equation}
A_1 v_1 = A_2 v_2
\end{equation}

Donde:
\begin{itemize}
    \item $A_1$ y $A_2$: Áreas de las secciones transversales en dos puntos diferentes de la tubería [m^2]
    \item $v_1$ y $v_2$: Velocidades del fluido en esos puntos [m/s]
\end{itemize}

Aplicando el Teorema de Bernoulli entre dos puntos (1 y 2) en un flujo de fluido, tenemos:
\begin{equation}
P_1 + \frac{1}{2} \rho v_1^2 + \rho g h_1 = P_2 + \frac{1}{2} \rho v_2^2 + \rho g h_2
\end{equation}

El Teorema de Bernoulli se puede aplicar en muchas situaciones prácticas, como en el diseño de alas de aviones, en sistemas de tuberías, en la medición de flujo con tubos Venturi, y en muchas otras aplicaciones donde el flujo de fluido es un factor importante.

Además de la ecuación principal de Bernoulli, es esencial entender los conceptos de presión estática, presión dinámica y presión total. La presión estática es la presión que se ejerce perpendicularmente a la superficie de un objeto en el fluido, la presión dinámica es la parte de la presión del fluido debida a su velocidad, y la presión total es la suma de la presión estática y la presión dinámica.

En resumen, el Teorema de Bernoulli proporciona una herramienta poderosa para analizar y entender el comportamiento de los fluidos en movimiento. A continuación, aplicaremos estos conceptos en varios ejercicios para reforzar el aprendizaje.


\actividad{Responder el siguiente cuestionario}
\begin{enumerate}
\item ¿Qué establece el Teorema de Bernoulli?
\item ¿Cómo se relacionan la velocidad y la presión de un fluido según el Teorema de Bernoulli?
\item ¿Qué condiciones deben cumplirse para aplicar el Teorema de Bernoulli?
\item ¿Qué representa cada término en la ecuación de Bernoulli?
\item ¿Cuál es la importancia de la ecuación de continuidad en el contexto del Teorema de Bernoulli?
\item ¿Cómo afecta la altura de un punto en un fluido en movimiento según el Teorema de Bernoulli?
\item ¿Qué es la presión dinámica y cómo se calcula?
\item ¿Qué aplicaciones prácticas tiene el Teorema de Bernoulli en la ingeniería?
\item ¿Cómo se mide la velocidad de un fluido utilizando un tubo Venturi?
\item ¿Cuál es la diferencia entre presión estática, presión dinámica y presión total?
\end{enumerate}
\actividad{Resolver los siguientes problemas}
\begin{enumerate}
\item En un tubo de Venturi, la sección 1 tiene un área de 0.1 m² y una velocidad de 2 m/s. La sección 2 tiene un área de 0.05 m². ¿Cuál es la velocidad del fluido en la sección 2?
\item Un fluido incompresible fluye a través de una tubería horizontal. En una sección donde la velocidad es 3 m/s, la presión es 150000 Pa. En otra sección, la velocidad es 6 m/s. ¿Cuál es la presión en esta segunda sección? (Suponga que la densidad del fluido es 1000 kg/m³).
\item En un sistema de tuberías, el agua fluye desde un punto A con una presión de 200000 Pa y una velocidad de 1 m/s a un punto B que está 5 m más alto. La velocidad del agua en el punto B es 2 m/s. ¿Cuál es la presión en el punto B?
\item Un avión vuela a una velocidad de 250 m/s a una altitud donde la densidad del aire es 0.9 kg/m³. ¿Cuál es la presión dinámica ejercida por el aire sobre el ala del avión?
\item En una tubería horizontal, la presión en la sección 1 es 180000 Pa y la velocidad es 1.5 m/s. En la sección 2, la presión es 120000 Pa. ¿Cuál es la velocidad en la sección 2? (Densidad del fluido: 1000 kg/m³).
\item Un fluido con densidad de 850 kg/m³ fluye a través de una tubería que se eleva 10 m. La velocidad en el punto inferior es 3 m/s y la presión es 210000 Pa. En el punto superior, la velocidad es 1.5 m/s. ¿Cuál es la presión en el punto superior?
\item En un tubo de Pitot, la velocidad del aire es medida en un avión en vuelo. Si la presión estática es 101325 Pa y la presión total medida es 105000 Pa, ¿cuál es la velocidad del aire? (Densidad del aire: 1.225 kg/m³).
\item Un líquido fluye a través de una tubería horizontal de 0.02 m² de área en la sección 1 y 0.01 m² en la sección 2. Si la velocidad en la sección 1 es 2 m/s, ¿cuál es la velocidad en la sección 2?
\item El agua fluye desde un tanque a través de un orificio a una velocidad de 5 m/s. Si la altura del agua sobre el orificio es 2 m, ¿cuál es la presión en el orificio? (Densidad del agua: 1000 kg/m³).
\item Un fluido con densidad de 950 kg/m³ fluye en una tubería donde la presión en un punto es 250000 Pa y la velocidad es 2.5 m/s. En otro punto, la presión es 200000 Pa y la altura es 3 m más alta. ¿Cuál es la velocidad en este segundo punto?
\end{enumerate}

\end{document}
